\documentclass[11pt]{article}
\usepackage[margin=1in]{geometry}
\usepackage{amsmath,amssymb,amsthm,hyperref}
\hypersetup{colorlinks=true,urlcolor=blue,linkcolor=blue}
\hypersetup{pdftitle={Mathematical Specification of the Certificate Software},
            pdfsubject={Algebra presentations, exact certificates, and reconstruction}}
\setlength{\parindent}{0pt}
\setlength{\parskip}{0.4em}
\newcommand{\Q}{\mathbb Q}
\newcommand{\Z}{\mathbb Z}
\newcommand{\A}{\mathcal A}
\newcommand{\ww}{\mathsf w}
\newtheorem{proposition}{Proposition}
\title{Mathematical Specification of the Certificate Software}
\author{}
\date{27 October 2026}
\begin{document}
\maketitle

\section{Purpose and boundary}
The installed package \texttt{zipper\_algebra} studies presentations motivated by zipper relations in quantum
topology and their untwisting quotients. Its initial maintained operations
construct explicit presentations, expand recorded proofs, compose two-sided
ideal-membership certificates, and check the resulting identities exactly.
These operations support research on the relations among sliders and on
quotients related to BMW algebras. They do not, by themselves, prove finite
generation, a dimension formula, termination of the geometric zipper algorithm,
or the topological origin of its relations.

This specification describes implemented algebraic operations. The full zipper
algorithm, thesis exposition and diagrams are separate
work. Historical experiment scripts remain available with their original
evidence, but their mathematical claims require individual verification.

\section{Ambient algebras and conventions}
Put $P=\Q[q,t]$ and $K=\Q(q,t)$, where $q,t$ are algebraically independent central
indeterminates. For $n\ge2$, use the free associative algebra
\[
 F_n=P\langle x_1,\ldots,x_{n-1},y_1,\ldots,y_{n-1}\rangle.
\]
A word is an ordered tuple of indices in that displayed generator order,
starting at one. The empty tuple represents the unit. Concatenation denotes
multiplication, in the order written. In the braid group algebra, $x_j$ maps to
$\sigma_j$ and $y_j$ to $\sigma_j^{-1}$.

The presentation explicitly includes the inverse relations $x_jy_j-1$ and
$y_jx_j-1$, adjacent braid relations and far commutations, together with
redundant inverse or mixed-sign versions used by the reduction engine. Let
$J_n$ be the ideal generated by the braid and inverse relations. Passing from
$F_n$ to $F_n/J_n$ identifies freely different words representing the same braid.
Consequently, an identity in the braid group algebra need not be a literal
free-algebra identity. The maintained verifier checks the latter. A certificate
using braid simplification must supply the corresponding braid/inverse terms
or use a separately specified group-algebra checker.

\section{The presentation builder}
The function \texttt{build\_relations\_ring} lives in \path{src/zipper_algebra/zipper_in_ms.py}.
Its front end is the class \texttt{Presentation} in \path{src/zipper_algebra/presentation.py}.
A presentation carries the strand
count, untwisting choice, framing and twist powers, level limits, zero-slider
relations and named additional relations. Its symbolic relation content is
hashed; readable relation labels alone do not identify an algebra.

\subsection{Sliders}
For a block of $i$ strands starting at $j$, define
\[
 U^+_{i,j}=x_jx_{j+1}\cdots x_{j+i-2},\qquad
 U^-_{i,j}=y_jy_{j+1}\cdots y_{j+i-2}.
\]
The implemented recursion is
\begin{align*}
 S_{2,j}&=q y_j+(1-q)-x_j,\\
 S_{i,j}&=(q^{i-1}U^-_{i,j}-U^+_{i,j})S_{i-1,j}\quad(3\le i\le n),
\end{align*}
with $1\le j\le n-i+1$. The unshifted slider $X_i$ means $S_{i,1}$.

\subsection{Zipper and canonical untwisting relations}
The level-two zipper relator is
\[
 Z_{2,j}=(q y_{j+1}+1-q-x_{j+1}-q y_jy_{j+1}+x_jx_{j+1})S_{2,j}.
\]
For $i\ge3$, the emitted higher zipper relator is
\[
 Z_{i,j}=\bigl(q^{i-1}U^-_{i,j+1}-U^+_{i,j+1}
      -q^{i-1}y_jU^-_{i,j+1}+x_jU^+_{i,j+1}\bigr)S_{i,j}.
\]
Here $i\le n-1$ and $1\le j\le n-i$. These formulas specify the current
builder; their equivalence to any alternate geometric formulation is a
separate mathematical obligation.

In the canonical untwisting presentation,
\[
 T_{2,j}=(x_j-t)S_{2,j},\qquad
 T_{i,j}=\bigl(x_{j+i-2}\cdots x_j
                 -q^{i-1}y_{j+i-2}\cdots y_j\bigr)S_{i,j}\quad(i\ge3).
\]
In the braid quotient, left multiplication by $U^+_{i,j}$ converts the second
equation to $(\Theta_{i,j}-q^{i-1})S_{i,j}=0$, where
$\Theta_{i,j}=U^+_{i,j}(x_{j+i-2}\cdots x_j)$. This equivalence uses invertibility;
the two expressions are different polynomials in the free algebra.

\subsection{Experimental parameter choices}
The builder also implements a level-two power relation $(x_j^a-t)S_{2,j}$,
the squared relation $(x_j-t)^2S_{2,j}$, and a historical option
$(x_j^2-t^2)S_{2,j}$. The twist-power parameter affects level three only:
for exponent $m>1$, the emitted relation is
$(\Theta_{3,j}^m-q^{2m})S_{3,j}$. At higher levels the canonical relation remains.
The software does not yet accept an arbitrary untwisting polynomial as input.
No structural theorem about a variant follows simply from choosing a flag.

The public presentation object rejects unknown untwisting choices, noninteger
strand counts and powers, level bounds outside $2,\ldots,n$, repeated additional
relation labels and targets requiring more strands than are present. Numeric
coefficients use either $\Q$ or $\Z/m\Z$. Calling the latter a field requires
the additional hypothesis that $m$ is prime; the coefficient backend does not
perform a primality proof. Nonunits cannot be inverted.

\section{Certificate data and exact conclusion}
Let $r_1,\ldots,r_s\in F_n$ be the ordered defining relations of a specified
presentation. A certificate for $f\in F_n$ is a finite list of terms
$(u_\nu,i_\nu,v_\nu,N_\nu,e_\nu)$, together with an ordered family of nonzero
polynomials $a_1,\ldots,a_h\in\Z[q,t]$. Here $u_\nu,v_\nu$ are words,
$i_\nu$ names an actual relation, $N_\nu\in P$, and
$e_\nu\in\mathbb N^h$. The intended identity is
\begin{equation}\label{eq:cert}
 f=\sum_\nu\frac{N_\nu}{D_\nu}u_\nu r_{i_\nu}v_\nu,
 \qquad D_\nu=\prod_{k=1}^h a_k^{e_{\nu k}},
\end{equation}
in $K\langle x_1,\ldots,y_{n-1}\rangle$. Atom irreducibility is unnecessary
for checking this identity; it is relevant to reconstruction efficiency.
The atom search bounds do not limit the exponents in a composed certificate.

Define $E_k=\max_\nu e_{\nu k}$, with zero maxima for an empty certificate,
and $M=\prod_k a_k^{E_k}$. Since $P$ is an integral domain and each atom is
nonzero, $M\ne0$. The exact checker compares the polynomials
\begin{equation}\label{eq:clear}
 Mf\quad\text{and}\quad
 \sum_\nu N_\nu\left(\prod_k a_k^{E_k-e_{\nu k}}\right)
                   u_\nu r_{i_\nu}v_\nu.
\end{equation}
Equality proves \eqref{eq:cert} and hence membership of $f$ in the two-sided
ideal generated by the relations over $K$. It also proves membership in the
corresponding localization of $P$, inverting the used atoms. A specialization
of the rational identity is available wherever its displayed denominators are
nonzero. The polynomial identity \eqref{eq:clear} remains meaningful at their
zeros, but need not imply the specialized membership statement there.

The certificate boundary rejects zero atoms, nonintegral atom coefficients,
negative or noninteger polynomial exponents, denominator vectors of the wrong
length, inexact numerical coefficients, out-of-range word indices and unknown
relation labels. Both verifiers validate the supplied presentation identity.
Objects are revalidated at the checking boundary because their data are mutable.
A malformed certificate raises a validation error. A well-formed false identity
returns a report with \texttt{pass=False}.

Legacy files without a presentation hash require an explicitly supplied
presentation. Their historical origin does not replace verification against
those actual relations. New certificates should retain the presentation hash
and explicit target. Named sliders, additional named targets and direct
polynomial targets share the same validated target resolution.

\section{Independent polynomial identity testing}
The second checker evaluates the cleared identity on a product grid of distinct
rational integers. For each parameter, its bound includes both the target
degree plus the clearing-factor degree and the degrees of every cofactor term.
If the bounds are $(d_q,d_t)$, a grid with $(d_q+1)(d_t+1)$ points suffices.

\begin{proposition}
If a polynomial in $\Q[q,t]$ of bidegree at most $(d_q,d_t)$ vanishes on such a
product grid, it is zero.
\end{proposition}
\begin{proof}
Fix a grid value of $t$ and apply the univariate root bound in $q$. Each
coefficient in $q$ then vanishes at all $d_t+1$ values of $t$, so the same
root bound makes each coefficient zero. Apply this to every free-word coefficient.
\end{proof}

The checker evaluates the polynomial product $M/D_\nu$ directly, rather than
dividing its evaluations. Thus grid points at poles cause no undefined
division. Negative numerator exponents are rejected: a Laurent polynomial
cannot use this polynomial degree bound without additional clearing.
This exact finite-grid argument differs from checking a few selected parameter
values. Its arithmetic is separate from the symbolic accumulation, but it
shares the presentation builder, input validation and target resolution.

\section{Proof expansion and contextual composition}
A basis proof is an acyclic directed graph expressing each basis element as a
sum of contextual multiples of input relations and other basis elements.
Forward references are allowed; cycles and missing dependencies are errors.
Substitution respects multiplication order. If
\[
 g_k=\sum_\nu c_{k\nu}L_{k\nu}r_{i_{k\nu}}R_{k\nu},
 \qquad f=\sum_s m_s u_sg_{k_s}v_s,
\]
composition produces terms with coefficient $m_sc_{k_s\nu}$ and contexts
\[
          (u_sL_{k_s\nu},\ i_{k_s\nu},\ R_{k_s\nu}v_s).
\]
No commutation or cancellation of free words is implicit. Repeated uses of one
basis element remain separate until their full contexts agree and coefficients
can be added. The number of multipliers must equal the number of skeleton steps,
and an explicit nonzero skeleton remainder precludes membership composition.

For sampled multipliers, each sample reduction must have the same ordered
sequence of $(u_s,k_s,v_s)$ as the anchor, where $k_s$ is the leading-word key
rather than the engine's changing numerical index. The data retain one
multiplier per occurrence. A sample with a different contextual schedule is
excluded. A successful fit is a candidate; the resulting identity still needs
exact verification.

\section{Exact reconstruction of candidates}
A sample consists of a distinct parameter point in $\Q^2$ and an exact rational
value. Repeated locations count once; conflicting values and floating-point
coefficients are rejected. For a fixed nonzero denominator proposal and a
numerator bidegree bound, interpolation uses the available fitting rows and
requires full column rank. A consistent rank-deficient system does not determine
a unique numerator and is not completed by arbitrarily setting free variables
to zero. Numerator coefficients belong to $\Q$, rather than necessarily $\Z$.

Univariate rational reconstruction fixes a monic denominator, solves the exact
linear system, cancels a polynomial greatest common divisor and normalizes the
result. The reduced denominator must be nonzero at every supplied point, and the
candidate must match every fitting and held-out value. This prevents a spurious
cross-multiplied equation $0=0$ from concealing an incorrect sample at a pole.

Held-out locations are reserved before fitting and denominator inference. Their
coefficients do not participate in inference of the factor family. A passing
held-out check is additional evidence, not a rational-function identity proof.
Line restrictions can cause distinct polynomial factors to coincide, so inferred
exponents are heuristic proposals. They are validated and have a general-search
fallback. For arbitrary declared factors, absence of a cancellable whole factor
does not imply coprimality of the numerator and denominator; stronger assertions
require irreducible, pairwise nonassociate factors.

Failed predictions, an unavailable common reduction shape and exhaustion of the
chosen degree or factor bounds do not show that more data cannot help. The
experimental driver retries within its sampling cap and reports an inconclusive
outcome when it lacks an exact certificate. In particular, reducing to a nonzero
remainder against an uncertified basis establishes no nonmembership conclusion,
even at that specialization. The historical exception name
\texttt{WrongClass} is retained as a subtype of \texttt{NeedMoreSamples}; it is
an inconclusive diagnostic. Candidate reconstruction is characteristic zero.

\section{Evidence and unsupported conclusions}
The historical large free-algebra certificate is retained for compatibility
checks. Smaller braid-algebra certificates use a different ambient equality and
must be described accordingly.

A finite computation, a small residual, agreement across specializations, or a
successful process exit does not certify a completed generic Gr\"obner basis.
A nonzero specialized remainder cannot alone establish generic nonmembership.
Dimension and trace claims require their own basis-completion and exact-arithmetic
evidence. Exact coefficient reconstruction is a maintained candidate-producing operation.
Solver execution, generic basis certification and the experiment drivers retain
separate completion requirements and remain experimental.

\section{Verification commands and source map}
The solver-free regressions and distribution checks cover strict certificate validation,
noncommuting contexts, repeated basis occurrences, schedule alignment, poles,
target degrees and historical certificate compatibility. See
\path{docs/TESTING.md} for the maintained commands and optional integrations.

The maintained implementation is under \path{src/zipper_algebra/}:
\begin{itemize}
\small
\setlength{\itemsep}{0pt}
\setlength{\parskip}{0pt}
\setlength{\parsep}{0pt}
\item Presentation and relation formulas: \path{presentation.py}, \path{zipper_in_ms.py}.
\item Certificate data: \path{cofactor/cofactor_io.py}, \path{cofactor/atoms.py}.
\item Symbolic and grid checks: \path{cofactor/verify.py}.
\item Proof graph and reduction: \path{cofactor/proof.py}, \path{cofactor/skeleton.py}.
\item Symbolic composition: \path{cofactor/compose.py}.
\item Coefficient reconstruction: \path{cofactor/lift.py}.
\end{itemize}
The experimental sampling driver is \path{tools/cofactor/pipeline.py}.
\end{document}
